\documentclass[11pt,a4paper]{article}
\usepackage[utf8]{inputenc}
\usepackage[T1]{fontenc}
\usepackage{lmodern}
\usepackage{geometry}
\usepackage{hyperref}
\usepackage{longtable}
\usepackage{array}
\usepackage{xcolor}
\geometry{margin=2.5cm}

\title{Documentation: \texttt{newsUpdate.ps1}}
\author{ScriptBox Project}
\date{\today}

\begin{document}
\maketitle
\tableofcontents
\newpage

\section{Purpose and Scope}
This document explains how the \texttt{src/newsUpdate.ps1} script works, how another user can set it up, and the current project status with TODOs and future plans.

\subsection{What the script does}
The script generates a personalized morning briefing by combining:
\begin{itemize}
  \item News from NewsAPI (tag-based queries),
  \item optional Hacker News results (via Algolia API),
  \item one deterministic daily recipe from TheMealDB,
  \item optional AI summary generation using a GitHub Copilot token.
\end{itemize}

It supports output in terminal, text file, and browser-rendered HTML.

\section{How the Script Works}
\subsection{Main execution flow}
At a high level, the script executes the following pipeline:
\begin{enumerate}
  \item Parse script parameters (API key, tags, output options, etc.).
  \item Validate minimal configuration (NewsAPI key placeholder check).
  \item Fetch one daily recipe from TheMealDB.
  \item For each configured tag:
    \begin{itemize}
      \item fetch NewsAPI articles,
      \item optionally fetch Hacker News items,
      \item normalize data into a shared structure,
      \item optionally generate an AI summary over collected headlines.
    \end{itemize}
  \item Print formatted output in terminal.
  \item Optionally save terminal output to a text file.
  \item Optionally build and open a responsive HTML page in a browser.
\end{enumerate}

\subsection{Important functions}
\begin{longtable}{>{\raggedright\arraybackslash}p{0.30\textwidth} >{\raggedright\arraybackslash}p{0.65\textwidth}}
\textbf{Function} & \textbf{Responsibility} \\
\hline
\texttt{Write-Color} & Color-aware terminal output with support for \texttt{-NoColor} and \texttt{-BrowserOnly}. \\
\texttt{HtmlEnc} & HTML-escapes strings before inserting data into generated HTML. \\
\texttt{Get-DailyRecipe} & Gets one recipe from TheMealDB based on day-of-year for deterministic daily rotation. \\
\texttt{Get-NewsEverything} / \texttt{Get-NewsTopHeadlines} & Calls NewsAPI endpoints and returns article arrays. \\
\texttt{Get-HackerNewsForTag} & Pulls stories from Algolia HN search API for each tag. \\
\texttt{Get-AiSummary} & Sends headlines to GitHub Copilot API and returns a short summary. \\
\texttt{Format-PublishedAt} & Converts ISO date strings to local time format \texttt{dd.MM.yyyy HH:mm}. \\
\texttt{Emit} & Centralized output helper used by terminal and file-save logic. \\
\end{longtable}

\subsection{Output modes}
\begin{itemize}
  \item \textbf{Terminal output} (default): formatted blocks by tag and source.
  \item \textbf{Text file output}: enabled with \texttt{-SaveToFile}; path configurable via \texttt{-OutputPath}.
  \item \textbf{Browser output}: enabled with \texttt{-OpenInBrowser} (plus terminal) or \texttt{-BrowserOnly} (no terminal output).
\end{itemize}

\section{Setup and User Configuration}
\subsection{Prerequisites}
\begin{itemize}
  \item Windows PowerShell (script currently written for PowerShell syntax used in \texttt{.ps1} files).
  \item Internet access.
  \item NewsAPI key from \url{https://newsapi.org/register}.
  \item Optional: GitHub token for AI summaries (Copilot-compatible endpoint).
\end{itemize}

\subsection{Recommended first-time setup}
\begin{enumerate}
  \item Copy repository and open a PowerShell terminal in the project root.
  \item Set your own API key values when running (recommended), instead of editing hardcoded defaults.
  \item Run the script once with a minimal command and verify data is returned.
\end{enumerate}

Example run:
\begin{verbatim}
powershell -ExecutionPolicy Bypass -File .\src\newsUpdate.ps1 `
  -NewsApiKey "<YOUR_NEWSAPI_KEY>" `
  -Tags "technology","AI","science" `
  -ArticlesPerTag 3
\end{verbatim}

Run with browser output and saved text file:
\begin{verbatim}
powershell -ExecutionPolicy Bypass -File .\src\newsUpdate.ps1 `
  -NewsApiKey "<YOUR_NEWSAPI_KEY>" `
  -OpenInBrowser -SaveToFile
\end{verbatim}

\subsection{Settings reference}
\begin{longtable}{>{\raggedright\arraybackslash}p{0.28\textwidth} >{\raggedright\arraybackslash}p{0.67\textwidth}}
\textbf{Parameter} & \textbf{Meaning / Typical Change} \\
\hline
\texttt{-NewsApiKey} & Required for NewsAPI calls. New users should always set their own value. \\
\texttt{-DisableHackerNews} & Disables HN contribution if only NewsAPI is desired. \\
\texttt{-GitHubToken} & Enables AI summary generation. Leave empty to disable summaries. \\
\texttt{-Tags} & Main personalization control; choose topics relevant to the user. \\
\texttt{-ArticlesPerTag} & Controls result count and output size. \\
\texttt{-Language} & NewsAPI language filter (validated list in script). \\
\texttt{-TopHeadlines} & Switch endpoint from \texttt{everything} to \texttt{top-headlines}. \\
\texttt{-NoColor} & Plain terminal output without ANSI colors. \\
\texttt{-SaveToFile} / \texttt{-OutputPath} & Persist output to text file for archival or automation. \\
\texttt{-OpenInBrowser} / \texttt{-BrowserOnly} & Generate and open HTML news dashboard. \\
\end{longtable}

\subsection{Sharing with another user}
For a clean handover, provide:
\begin{itemize}
  \item this documentation file,
  \item a short run command template,
  \item clear note that API keys/tokens must be user-owned and not committed.
\end{itemize}

\subsection{Quick start (60 seconds)}
\begin{enumerate}
  \item Open PowerShell in project root.
  \item Run with your own NewsAPI key.
  \item Confirm that at least one tag prints article entries.
\end{enumerate}

\begin{verbatim}
powershell -ExecutionPolicy Bypass -File .\src\newsUpdate.ps1 `
  -NewsApiKey "<YOUR_NEWSAPI_KEY>"
\end{verbatim}

Expected result: a timestamped briefing with recipe block and one section per tag.

\subsection{Troubleshooting}
\begin{itemize}
  \item \textbf{No output or script blocked:} check PowerShell execution policy and run with \texttt{-ExecutionPolicy Bypass}.
  \item \textbf{NewsAPI errors:} verify key validity, quota, and endpoint limits on your plan.
  \item \textbf{Empty article sections:} try broader tags (for example \texttt{technology}) and set \texttt{-TopHeadlines} off.
  \item \textbf{AI summary missing:} ensure \texttt{-GitHubToken} is set and token has correct access.
  \item \textbf{Browser did not open:} run with \texttt{-OpenInBrowser} and verify default browser association on the OS.
\end{itemize}

\section{Current State, TODOs, and Future Plans}
\subsection{Current state (as implemented)}
\begin{itemize}
  \item Multi-source aggregation works (NewsAPI + optional HN + recipe).
  \item Optional AI summary integration is present.
  \item Output supports terminal, file, and browser dashboard.
  \item Error handling exists around external API calls.
\end{itemize}

\subsection{TODOs (recommended next tasks)}
\begin{itemize}
  \item Remove hardcoded default API key from source and require secure input (env var or argument only).
  \item Add per-source request timeout and retry/backoff strategy.
  \item Add duplicate article detection between NewsAPI and HN results.
  \item Add automated tests for formatter logic and date conversion.
  \item Add logging mode for scheduled/background runs.
\end{itemize}

\subsection{Future plans}
\begin{itemize}
  \item Optional categories/profiles (e.g., \texttt{work}, \texttt{gaming}, \texttt{finance}).
  \item Caching layer to reduce repeated API calls within short intervals.
  \item Export options such as Markdown/JSON for downstream processing.
  \item Config file support (e.g., \texttt{newsUpdate.config.json}) for non-technical users.
  \item Optional notification integration (email/Teams/Slack).
\end{itemize}

\subsection{Known limitations}
\begin{itemize}
  \item News relevance depends on external provider indexing and API quotas.
  \item No built-in deduplication currently; similar stories can appear from multiple sources.
  \item Secrets are parameter-based and not yet managed through a dedicated secret store.
  \item HTML output path is temporary by default and is not retained unless separately archived.
\end{itemize}

\section{Maintenance Notes}
\begin{itemize}
  \item Keep endpoint URLs and model names reviewed periodically; provider APIs can evolve.
  \item Rotate tokens/keys and avoid storing secrets in version control.
  \item Re-check CSS and layout in generated HTML on mobile after major edits.
\end{itemize}

\end{document}

